\documentclass[12pt]{article}
\usepackage[margin=1in]{geometry}
\usepackage{graphicx}
\usepackage{booktabs}
\usepackage{amsmath}

\title{Replicating Table 1, Panel B from\\Chernozhukov et al.\ (2017)}
\author{Aaron Chotzen-Jenner}
\date{\today}

\begin{document}
\maketitle

\section{The Original Paper}

Chernozhukov, Chetverikov, Demirer, Duflo, Hansen, and Newey (2017),
``Double/Debiased/Neyman Machine Learning of Treatment Effects,''
\textit{American Economic Review: Papers \& Proceedings}, 107(5): 261--265,
introduces the Double/Debiased Machine Learning (Double ML) estimator for
causal inference in settings where high-dimensional nuisance parameters must
be estimated. The central research question is how to obtain valid,
root-$N$ consistent estimates of a treatment effect parameter $\theta_0$
when the confounding functions are estimated using flexible machine learning
methods that introduce regularization bias.

The empirical strategy rests on two ingredients. First, the estimating
equation is made Neyman-orthogonal, meaning it is locally insensitive to
small errors in the nuisance parameter estimates. Second, cross-fitting
is used to split the sample into folds, training the nuisance functions
on held-out data to avoid overfitting bias. Together these two steps
remove the regularization bias that would otherwise contaminate a naive
plug-in estimator.

\section{The Result of Interest}

The specific result I replicate is the Mean Average Treatment Effect (ATE)
reported in Table 1, Panel B, row ``ATE (2 fold),'' Ensemble column of
the online appendix. This result corresponds to the partially linear model
with 2-fold cross-fitting. The paper reports a Mean ATE of \$9,043 with
a standard error of 1,432.

This result is the headline finding of the empirical illustration: after
flexibly controlling for income and other confounders using an ensemble of
machine learning methods, 401(k) eligibility raises net financial assets
by roughly \$9,000. It is the cleanest single number to target because
it corresponds to the default specification in the replication code
(\texttt{Master.R}), uses the partially linear model, and focuses on the
Ensemble method, which combines the machine learning approaches.

\section{Data and Methods}

The data come from the 1991 Survey of Income and Program Participation
(SIPP), the same dataset used in Chernozhukov and Hansen (2004). The
sample contains 9,915 observations. The outcome variable \texttt{net\_tfa}
is net total financial assets, measured in dollars. The treatment variable
\texttt{e401} is a binary indicator for eligibility to enroll in a 401(k)
plan. The nine covariates used in the model are age, income, years of
education, family size, and binary indicators for married status,
two-earner household, defined benefit pension, IRA participation, and
home ownership.

\texttt{preprocess.R} loads the raw Stata file \texttt{sipp1991.dta} from
\texttt{input/}, selects the eleven relevant variables, and writes a clean
CSV to \texttt{temp/clean\_data.csv}. \texttt{analysis.R} loads that file,
sources the original helper functions \texttt{ML\_Functions.R} and
\texttt{Moment\_Functions.R} from \texttt{input/}, and calls
\texttt{DoubleML()} with the partially linear model specification and
2-fold cross-fitting. To reduce computation time, the estimation is
repeated over 10 random sample splits rather than the 100 splits used in
the original paper. Parallel computation across 4 cores is used to
accelerate estimation. The Mean ATE, Median ATE, and their standard errors
are computed following the paper's formulas exactly. Results are written
to \texttt{output/tables/main\_result.tex}, which is
\texttt{\textbackslash input{}}ed directly into this document.

\section{Replication Results}

\input{../output/tables/main\_result.tex}

\input{../output/tables/comparison\_table.tex}

The replication estimate reported in Table 1 closely matches the
corresponding result in Chernozhukov et al.\ (2017). The estimated average
treatment effect of 401(k) eligibility on net financial assets is very
similar to the published estimate, and the associated standard errors are
nearly identical. Given that the original study averages over 100 random
sample splits while this replication uses only 10, small differences of
this magnitude are expected. Overall, the replication successfully
reproduces the paper's central empirical finding that 401(k) eligibility
substantially increases net financial assets.

\section{What I Learned}

Once the helper functions were sourced correctly, reproducing the estimation
was straightforward. The hardest part was path management: the original code
assumed all files lived in the same working directory, while the
Gentzkow-Shapiro structure separates \texttt{input/}, \texttt{code/}, and
\texttt{temp/}. In particular, \texttt{Moment\_Functions.R} internally
sources \texttt{ML\_Functions.R} with a bare filename, which required
patching to use an absolute path. This is a common reproducibility issue:
code that works in the author's environment can break when placed in a
different directory structure.

More broadly, this replication taught me that Double ML is remarkably
stable: six very different machine learning methods (RLasso, Trees, Random
Forest, Boosting, Neural Net, Ensemble) all produce estimates between
approximately \$7,900 and \$9,200, with substantial overlap in their confidence
intervals. This stability is exactly what the Neyman orthogonality theory
predicts, and seeing it empirically is more convincing than any theorem.
If I were the original authors, the one thing that would have made
replication easier is absolute rather than relative paths in the helper
files, along with a brief README note specifying how the scripts should be
run.

\begin{thebibliography}{9}
\bibitem{paper}
Chernozhukov, V., Chetverikov, D., Demirer, M., Duflo, E., Hansen, C.,
and Newey, W. (2017). ``Double/Debiased/Neyman Machine Learning of
Treatment Effects.'' \textit{American Economic Review: Papers \&
Proceedings}, 107(5): 261--265.
\end{thebibliography}

\end{document}
